\documentclass{article}
\usepackage{graphicx}
\usepackage[style=ieee]{biblatex} % Set the reference style to IEEE
\usepackage{xcolor}
\usepackage{hyperref}
\usepackage{titletoc}

\definecolor{mygreen}{RGB}{0,128,0}

\usepackage{array} % For customizing the table
\usepackage{booktabs} % For better quality horizontal lines
\usepackage{graphicx} % Package for including images
\usepackage{float}

% Define margins
\usepackage[margin=1in]{geometry}

\renewcommand{\contentsname}{\textcolor{mygreen}{Table of Contents}}

\begin{document}

\begin{titlepage}
    \centering
    % University logo
    \includegraphics[width=0.24\textwidth]{logo.jpg}
    \par\vspace{2cm}

    % University name and course details
    {\Large \textbf{Your Institution Name} \par}
    \vspace{0.5cm}
    {\large Major \par}
    {\large Code - Subject \par}
    \vspace{3cm}

    % Lab and activity details
    {\large \textbf{Topic} \par}
    {\large Topic Detailed \par}
    {\large Specific topic \par}
    \vspace{3cm}

    % Team members
    {\large \textbf{Team} \par}
    \vspace{0.5cm}
    \begin{tabular}{ll}
    Name1 & name1@name1.com \\
    Name1 & name1@name1.com \\
    Name1 & name1@name1.com \\
    Name1 & name1@name1.com \\
    Name1 & name1@name1.com \\
    Name1 & name1@name1.com \\
    \end{tabular}
    \par\vspace{3cm}

    % Date
    {\large Month day year \par}
\end{titlepage}

\tableofcontents % Insert table of contents

\newpage % Page break to separate table of contents from document content

% Document content
\section{Introduction}\label{sec:example}
Here begins the content of the article in section \ref{sec:example}, and there's also a reference [\ref{ref:wellUniv}].

\renewcommand{\tablename}{Table}
\begin{table}[htbp]
  \centering
  \caption{Example of a 4x3 table with marked outline}
  \begin{tabular}{|c|c|c|c|} % 4 columns, all centered and with marked outline
    \hline
    Header 1 & Header 2 & Header 3 & Header 4 \\
    \hline
    Cell 1 & Cell 2 & Cell 3 & Cell 4 \\
    \hline
    Cell 5 & Cell 6 & Cell 7 & Cell 8 \\
    \hline
    Cell 9 & Cell 10 & Cell 11 & Cell 12 \\
    \hline
  \end{tabular}
   \label{tab:Test}
\end{table}

There's also a table \ref{tab:Test}.

\renewcommand{\tablename}{Figure}
\begin{figure}[htbp]
  \centering
  \includegraphics[width=0.3\textwidth]{logo.jpg}
  %\includegraphics[scale=0.2]{university_logo.png}
  \caption{Description of the figure (Image designed by freepik).}
  \label{fig:example}
\end{figure}

And there's a figure \ref{fig:example}.\\

Finally, a list:

\begin{itemize}
  \item Item 1
  \item Item 2
  \item Item 3\\\\
\end{itemize}


Use the following page to document references:
\href{https://www.citethisforme.com/}{https://www.citethisforme.com/}

%\renewcommand{\refname}{Referencias}
\begin{thebibliography}{9}

\bibitem{ref} \label{ref:wellUniv}University Welfare Directorate, “National University Welfare Directorate,” visit counter. [Online]. Available: \href{https://welfare.edu.co/index.php}{https://welfare.edu.co/index.php}. [Accessed: March 5, 2023].  

\end{thebibliography}

\end{document}
